\documentclass[11pt,a4paper]{article}
\usepackage[utf8]{inputenc}
\usepackage[margin=1in,headheight=14pt]{geometry}
\usepackage{xcolor}
\usepackage{listings}
\usepackage{tcolorbox}
\usepackage{enumitem}
\usepackage{hyperref}
\usepackage{amsmath,amssymb}
\usepackage{fancyhdr}
\usepackage{microtype}
\usepackage{booktabs}
\usepackage{array}

% Hyperlink configuration
\hypersetup{
    colorlinks=true,
    linkcolor=blue!70!black,
    urlcolor=blue!70!black,
    citecolor=blue!70!black
}

% Colors for code listing
\definecolor{codegreen}{rgb}{0,0.6,0}
\definecolor{codegray}{rgb}{0.5,0.5,0.5}
\definecolor{codepurple}{rgb}{0.58,0,0.82}
\definecolor{backcolour}{rgb}{0.97,0.97,0.98}
\definecolor{termback}{rgb}{0.12,0.12,0.12}
\definecolor{termfore}{rgb}{0.9,0.9,0.9}

% Listings style for C code
\lstdefinestyle{cstyle}{
    backgroundcolor=\color{backcolour},   
    commentstyle=\color{codegreen}\itshape,
    keywordstyle=\color{blue}\bfseries,
    numberstyle=\tiny\color{codegray},
    stringstyle=\color{codepurple},
    basicstyle=\ttfamily\footnotesize,
    breakatwhitespace=false,         
    breaklines=true,                 
    captionpos=b,                    
    keepspaces=true,                 
    numbers=left,                    
    numbersep=7pt,                  
    showspaces=false,                
    showstringspaces=false,
    showtabs=false,                  
    tabsize=4,
    frame=single,
    rulecolor=\color{black!20},
    language=C
}

% Listings style for bash commands / scripts
\lstdefinestyle{bashstyle}{
    backgroundcolor=\color{backcolour},
    commentstyle=\color{codegreen}\itshape,
    keywordstyle=\color{blue}\bfseries,
    numberstyle=\tiny\color{codegray},
    stringstyle=\color{codepurple},
    basicstyle=\ttfamily\footnotesize,
    breaklines=true,
    frame=single,
    rulecolor=\color{black!20},
    numbers=none,
    showstringspaces=false,
    language=bash
}

% Listings style for terminal output
\lstdefinestyle{termstyle}{
    backgroundcolor=\color{backcolour},
    basicstyle=\ttfamily\footnotesize,
    breaklines=true,
    frame=single,
    rulecolor=\color{black!20},
    numbers=none,
    showstringspaces=false
}

\lstset{style=cstyle}

\pagestyle{fancy}
\fancyhf{}
\lhead{\small\textbf{CPE333 Operating Systems (1/2026)}}
\rhead{\small PS3: Process Manipulation \& Monitoring}
\cfoot{\thepage}

\begin{document}
\thispagestyle{plain}

\begin{center}
    {\LARGE \textbf{CPE333 Operating Systems (1/2026)}}\\[4pt]
    {\Large \textbf{Problem Session 3: Process Manipulation and Monitoring}}\\[8pt]
    \rule{\textwidth}{0.8pt}
\end{center}

\vspace{-0.5em}
\noindent
\begin{tabular*}{\textwidth}{@{\extracolsep{\fill}}ll}
    \textbf{1.} Name: Thanaboon Tikaew & \textbf{ID:} 67070501021 \\[3pt]
    \textbf{2.} Name: Chanon Lhumsa-ard & \textbf{ID:} 67070501059 \\[3pt]
    \textbf{3.} Name: Theeraphat Jaingam & \textbf{ID:} 67070501063 \\[3pt]
    \textbf{4.} Name: Siriwan Yindeephot & \textbf{ID:} 67070501075 \\[4pt]
\end{tabular*}

\vspace{2pt}
\noindent\textbf{Environment:} Ubuntu 24.04 on WSL2, gcc 15.2.0, 8 CPUs\\[-4pt]
\rule{\textwidth}{0.8pt}

\vspace{0.5em}

\section{Task 1: Important Linux Commands for Process Manipulation}

\subsection{Difference between \texttt{top} and \texttt{ps}}

Both \texttt{top} and \texttt{ps} are used to inspect active processes on a Linux system, but they serve different operational purposes.

\begin{table}[h!]
\centering
\small
\begin{tabular}{@{}p{0.46\textwidth}p{0.48\textwidth}@{}}
\toprule
\textbf{\texttt{ps}} & \textbf{\texttt{top}} \\
\midrule
\textbf{Static snapshot:} Prints a point-in-time snapshot of processes and exits immediately. & 
\textbf{Dynamic monitor:} Continuously refreshes and displays real-time system activity. \\
\addlinespace
\textbf{Batch / Scripting friendly:} Ideal for scripting, pipelines, and grep (e.g., \texttt{ps aux | grep ssh}). & 
\textbf{Interactive:} Stays active until user exits (typically by pressing \texttt{q}). \\
\addlinespace
\textbf{Complete listing:} Outputs all matching processes in standard output stream. & 
\textbf{Dashboard view:} Summarizes top consumers fitting the current terminal dimensions. \\
\addlinespace
\textbf{Non-sorting by default:} Requires specific CLI flags to sort by resource metrics. & 
\textbf{Sorted automatically:} Sorts by CPU or memory usage by default. \\
\addlinespace
\textbf{Zero persistent overhead:} Exits immediately after reading \texttt{/proc}. & 
\textbf{Continuous overhead:} Consumes minor CPU while periodically refreshing. \\
\bottomrule
\end{tabular}
\caption{Comparison between \texttt{ps} and \texttt{top}.}
\end{table}

\noindent\textbf{Usage Example:}
\begin{lstlisting}[style=bashstyle]
$ ps aux                      # one-time process snapshot
$ top                         # live, interactive monitoring dashboard (quit with 'q')
\end{lstlisting}

\subsection{Purpose of the \texttt{nice} command}

The \texttt{nice} command launches a program with a modified \emph{scheduling priority}. Each process is assigned a ``niceness'' value determining how the kernel scheduler prioritizes CPU time allocations. A higher niceness represents a lower priority (the process is ``nicer'' to other tasks), whereas a lower (or negative) niceness elevates execution priority.

\begin{itemize}[leftmargin=1.5em]
    \item \textbf{Niceness range:} From $-20$ (highest priority) to $+19$ (lowest priority).
    \item \textbf{Default niceness:} $0$.
    \item \textbf{Privileges:} Only root/privileged users can assign negative niceness values.
\end{itemize}

\begin{lstlisting}[style=bashstyle]
$ nice -n 10 ./my_program      # run with niceness 10 (lower priority)
$ sudo nice -n -5 ./my_program # run with niceness -5 (higher priority, root required)
$ renice -n 5 -p 1234          # adjust priority of running process PID 1234
\end{lstlisting}

\subsection{How to kill processes and jobs}

The \texttt{kill} command delivers signals to processes identified by their Process ID (PID) or Job ID. The most frequently used signals include \texttt{SIGTERM} (15) for graceful termination and \texttt{SIGKILL} (9) for immediate unconditional termination.

\begin{lstlisting}[style=bashstyle]
$ ps aux | grep my_program      # identify target process PID
$ kill 1234                     # send SIGTERM (graceful shutdown)
$ kill -9 1234                  # send SIGKILL (forceful immediate termination)
$ killall my_program            # terminate all processes with matching name
$ pkill my_program              # terminate processes matching regex pattern
$ jobs                          # view current shell background/stopped jobs
$ kill %1                       # terminate job 1 from job table
\end{lstlisting}

\begin{lstlisting}[style=termstyle]
$ ps aux | grep ss1_1
ssiriwan  5678  0.0  0.1  12345  2345 pts/0  S  10:00  0:00 ./ss1_1.sh
$ kill 5678                     # gracefully terminate script process
$ kill -9 5678                  # forcibly terminate process if SIGTERM is blocked
\end{lstlisting}

\newpage

\section{Task 2: Running Shell Scripts in Foreground and Background}

\subsection{First situation: \texttt{ss1\_1.sh}}

The script \texttt{ss1\_1.sh} performs a 10-second sleep loop:

\begin{lstlisting}[style=bashstyle]
#!/bin/bash
for i in $(seq 1 10); do
  sleep 1
done
\end{lstlisting}

\subsubsection*{Running in Foreground (\texttt{./ss1\_1.sh})}
\begin{itemize}[leftmargin=1.5em]
    \item The process occupies the terminal standard input/output streams.
    \item The shell prompt is blocked until the script completes (10 seconds).
    \item The process can be interrupted directly via \texttt{Ctrl+C} (\texttt{SIGINT}).
\end{itemize}

\begin{lstlisting}[style=termstyle]
$ ./ss1_1.sh        # prompt blocked for ~10s, then returns
\end{lstlisting}

\subsubsection*{Running in Background (\texttt{./ss1\_1.sh \&})}
\begin{itemize}[leftmargin=1.5em]
    \item The trailing \texttt{\&} tells the shell to fork the process without blocking the parent shell.
    \item The shell outputs the job ID and PID (e.g., \texttt{[1] 5678}) and prompt returns immediately.
    \item The user can execute further commands concurrently while the script runs in the background.
\end{itemize}

\begin{lstlisting}[style=termstyle]
$ ./ss1_1.sh &
[1] 5678            # job number [1], PID 5678, prompt returns immediately
\end{lstlisting}

\subsubsection*{Discussion}
\begin{itemize}[leftmargin=1.5em]
    \item \textbf{Foreground execution:} The shell performs a \texttt{wait()} system call for the child process, preventing interactive user input until termination.
    \item \textbf{Background execution:} The shell creates the child process but does not block on \texttt{wait()}, returning control to the terminal input loop immediately.
\end{itemize}

\subsection{Second situation: \texttt{ss1\_2.sh}}

The script \texttt{ss1\_2.sh} sleeps for 1000 seconds:

\begin{lstlisting}[style=bashstyle]
#!/bin/bash
sleep 1000
\end{lstlisting}

\subsubsection*{Suspending Process with \texttt{Ctrl+Z}}
Pressing \texttt{Ctrl+Z} delivers the \texttt{SIGTSTP} signal to the foreground process group. The process is paused (state \texttt{T} / Stopped) and added to the shell's job table without being killed.

\begin{lstlisting}[style=termstyle]
$ ./ss1_2.sh
^Z
[1]+  Stopped                 ./ss1_2.sh
$ jobs
[1]+  Stopped                 ./ss1_2.sh
\end{lstlisting}

\subsubsection*{Resuming with \texttt{fg} and \texttt{bg}}
\begin{itemize}[leftmargin=1.5em]
    \item \textbf{\texttt{fg \%1}:} Resumes the suspended job in the \textbf{foreground} by sending \texttt{SIGCONT}. The terminal prompt blocks until completion.
    \item \textbf{\texttt{bg \%1}:} Resumes the suspended job in the \textbf{background} by sending \texttt{SIGCONT}. The prompt remains available for further commands.
\end{itemize}

\begin{lstlisting}[style=bashstyle]
$ fg %1     # resume job 1 in foreground
$ bg %1     # resume job 1 in background
\end{lstlisting}

\newpage

\section{Task 5: STCF (Shortest Time to Completion First) Scheduler}

\subsection{Algorithm Overview}

Shortest Time to Completion First (STCF), also referred to as \emph{Shortest Remaining Time First (SRTF)}, is a \textbf{preemptive} CPU scheduling policy. At every discrete time unit, the scheduler evaluates all available processes that have arrived and allocates the CPU to the process with the shortest remaining burst time.

\begin{itemize}[leftmargin=1.5em]
    \item \textbf{Zero-overhead model:} Scheduler evaluation and context switching take zero CPU time.
    \item \textbf{Arrival boundary:} A process arriving at time $T$ is available at time $T-1$'s completion, ready for evaluation at time $T$.
    \item \textbf{Tie-breaking rule 1:} If multiple processes arrive simultaneously, prioritize by lowest \texttt{process\_id}.
    \item \textbf{Tie-breaking rule 2:} If multiple active processes share identical remaining burst time, prioritize by lowest \texttt{process\_id}.
\end{itemize}

\subsection{Implementation}

\begin{lstlisting}[language=C]
/*
 * Scheduler function for PS3.c (STCF / Shortest Remaining Time First)
 * processes[] : array of processes (pid, arrival, burst)
 * n           : number of processes
 */
void stcf(struct Process p[], int n) {
    int remaining[MAX];
    int completed = 0;
    int t = 0;
    /* Initialize remaining burst times */
    for (int i = 0; i < n; i++)
        remaining[i] = p[i].burst;
    while (completed < n) {
        int chosen = -1;
        /* Select available process with minimum remaining burst time;
           resolve ties using lowest pid */
        for (int i = 0; i < n; i++) {
            if (p[i].arrival <= t && remaining[i] > 0) {
                if (chosen == -1 ||
                    remaining[i] < remaining[chosen] ||
                    (remaining[i] == remaining[chosen] && p[i].pid < p[chosen].pid)) {
                    chosen = i;
                }
            }
        }
        if (chosen == -1) {
            t++; /* CPU idle: advance time */
            continue;
        }
        /* Execute chosen process for 1 time slice */
        printf("Time %2d-%2d : Process %d\n", t, t + 1, p[chosen].pid);
        remaining[chosen]--;
        if (remaining[chosen] == 0)
            completed++;
        t++;
    }
}
\end{lstlisting}

\newpage

\subsection{Test Case Results}

\subsubsection{Case 1 (\texttt{case1.csv})}

\begin{center}
\begin{minipage}[t]{0.45\textwidth}
\centering
\textbf{Input Specification}
\vspace{0.3em}

\begin{tabular}{ccc}
\toprule
\textbf{PID} & \textbf{Arrival} & \textbf{Burst} \\
\midrule
1 & 0 & 7 \\
2 & 2 & 4 \\
3 & 4 & 1 \\
\bottomrule
\end{tabular}
\end{minipage}
\hfill
\begin{minipage}[t]{0.45\textwidth}
\centering
\textbf{Output Timeline}
\vspace{0.3em}

\begin{tabular}{cc}
\toprule
\textbf{Time Slot} & \textbf{Running Process} \\
\midrule
0--1   & Process 1 \\
1--2   & Process 1 \\
2--3   & Process 2 \\
3--4   & Process 2 \\
4--5   & Process 3 \\
5--6   & Process 2 \\
6--7   & Process 2 \\
7--8   & Process 1 \\
8--9   & Process 1 \\
9--10  & Process 1 \\
10--11 & Process 1 \\
11--12 & Process 1 \\
\bottomrule
\end{tabular}
\end{minipage}
\end{center}

\subsubsection{Case 2 (\texttt{case2.csv})}

\begin{center}
\begin{minipage}[t]{0.45\textwidth}
\centering
\textbf{Input Specification}
\vspace{0.3em}

\begin{tabular}{ccc}
\toprule
\textbf{PID} & \textbf{Arrival} & \textbf{Burst} \\
\midrule
1 & 0 & 2 \\
2 & 5 & 3 \\
3 & 5 & 3 \\
4 & 6 & 1 \\
\bottomrule
\end{tabular}
\end{minipage}
\hfill
\begin{minipage}[t]{0.45\textwidth}
\centering
\textbf{Output Timeline (CPU Idle: 2--5)}
\vspace{0.3em}

\begin{tabular}{cc}
\toprule
\textbf{Time Slot} & \textbf{Running Process} \\
\midrule
0--1   & Process 1 \\
1--2   & Process 1 \\
2--5   & \textit{(CPU Idle)} \\
5--6   & Process 2 \\
6--7   & Process 4 \\
7--8   & Process 2 \\
8--9   & Process 2 \\
9--10  & Process 3 \\
10--11 & Process 3 \\
11--12 & Process 3 \\
\bottomrule
\end{tabular}
\end{minipage}
\end{center}

\subsubsection{Case 3 (\texttt{case3.csv})}

\begin{center}
\begin{minipage}[t]{0.45\textwidth}
\centering
\textbf{Input Specification}
\vspace{0.3em}

\begin{tabular}{ccc}
\toprule
\textbf{PID} & \textbf{Arrival} & \textbf{Burst} \\
\midrule
2 & 0 & 5 \\
1 & 3 & 2 \\
3 & 6 & 1 \\
\bottomrule
\end{tabular}
\end{minipage}
\hfill
\begin{minipage}[t]{0.45\textwidth}
\centering
\textbf{Output Timeline}
\vspace{0.3em}

\begin{tabular}{cc}
\toprule
\textbf{Time Slot} & \textbf{Running Process} \\
\midrule
0--1  & Process 2 \\
1--2  & Process 2 \\
2--3  & Process 2 \\
3--4  & Process 1 \\
4--5  & Process 1 \\
5--6  & Process 2 \\
6--7  & Process 2 \\
7--8  & Process 3 \\
\bottomrule
\end{tabular}
\end{minipage}
\end{center}

\subsection{Discussion}

\begin{itemize}[leftmargin=1.5em]
    \item \textbf{Case 1 (Preemption Flow and Short-Job Prioritization):} Process 1 executes from $t = 0\text{--}2$ (remaining burst: $5$), whereupon Process 2 arrives with burst $4$ and immediately preempts Process 1 ($4 < 5$), followed by Process 3 arriving at $t = 4$ (burst $1$) and preempting Process 2 ($1 < 2$), which demonstrates how STCF enforces dynamic multi-level preemption to minimize response and waiting times for newly arriving short tasks before allowing preempted processes (Process 2, then Process 1) to resume and finish at $t = 7$ and $t = 12$ respectively.
    \item \textbf{Case 2 (CPU Idle State and Simultaneous Arrival Resolution):} After Process 1 completes at $t = 2$, the CPU remains idle from $t = 2\text{--}5$ until Process 2 and Process 3 arrive simultaneously with identical bursts of $3$, where Process 2 is selected first via the lower-PID tie-breaking rule ($\text{PID } 2 < \text{PID } 3$) and runs until $t = 6$ before being preempted by Process 4 (burst $1 < 2$), demonstrating that the scheduler robustly handles non-contiguous arrival gaps while applying deterministic tie-breaking to schedule simultaneous jobs before completing Process 2 ($t = 7\text{--}9$) and Process 3 ($t = 9\text{--}12$).
    \item \textbf{Case 3 (Tie-Breaking Behavior on Equal Remaining Time):} Process 2 executes from $t = 0\text{--}3$ (remaining: $2$) until Process 1 arrives with burst $2$, triggering a preemption because Process 1 has the lower PID ($\text{PID } 1 < \text{PID } 2$), yet when Process 3 arrives at $t = 6$ (burst $1$) while the resumed Process 2 also has $1$ remaining unit, Process 2 is allowed to continue running without preemption ($\text{PID } 2 < \text{PID } 3$), which proves that the PID tie-breaking rule functions symmetrically to both preempt a running process and protect an ongoing execution whenever candidate processes share identical remaining execution times.
    \item \textbf{Summary:} Overall, the STCF scheduler guarantees optimal average turnaround time by continuously dispatching the job closest to completion, seamlessly idling during arrival gaps, and resolving all collision scenarios through consistent PID-based tie-breaking.
\end{itemize}

\end{document}